%% kapa: exam=1 round=34 year=2023 period=2 form=A subject=appraisal_law subject_ko=감정평가관계법규 seq=39 number=39
%% source: Q-Net 감정평가사 기출문제 (cst00309) artlSeq=5211789 file=raw/1008/5211789/2023년도 제34회 감정평가사 1차시험 2교시.pdf
%% conversion: kordoc 4.24.51 -> tools/gen1.py (auto); review=proofread-sonnet; flags=none
\begin{question}{39}
도시 및 주거환경정비법령상 비용의 부담 등에 관한 설명으로 옳지 않은 것은?
\begin{choices}
  \item 사업시행자는 토지등소유자로부터 정비사업비용과 정비사업의 시행과정에서 발생한 수입의 차액을 부과금으로 부과·징수할 수 있다.
  \item 체납된 부과금 또는 연체료의 부과·징수를 위탁받은 시장·군수등은 지방세 체납처분의 예에 따라 부과·징수할 수 있다.
  \item 공동구점용예정자가 부담할 공동구의 설치에 드는 비용의 부담비율은 공동구의 점용예정면적비율에 따른다.
  \item 부담금의 납부통지를 받은 공동구점용예정자는 공동구의 설치공사가 착수되기 전에 부담금액의 3분의 1 이상을 납부하여야 한다.
  \item 시장·군수등은 시장·군수등이 아닌 사업시행자가 시행하는 정비사업의 정비계획에 따라 설치되는 임시거주시설에 대해서는 그 건설비용의 전부를 부담하여야 한다.
\end{choices}
\end{question}
\answer{5}
